\chapter{Модель угроз}

Модель угроз есть договор между протоколом и его пользователем: при каких допущениях что гарантируется. Всё, что не перечислено как гарантия, не гарантируется, и это сказано явно.

\section{Участники и допущения}

\begin{center}
\small
\begin{tabular}{p{3.2cm} p{10.4cm}}
\toprule
Участник & Допущение \\
\midrule
Рабочая нода & не доверена; может быть злонамеренной, скомпрометированной или отравленной инжектом через содержимое задачи, репозитория или зависимостей \\
Кор-нода, оркестратор & L0: доверены; L1: не более меньшинства скомпрометированных ключей при честном секвенсере; L2: не более $f$ злонамеренных из $3f+1$ \\
Среда & доверена в отношении целостности на всех уровнях, если не аттестована аппаратно; протокол защищает от кор-нод и агентов, но не от владельца среды \\
Провайдер модели & видит содержимое санбокса; доверен в отношении конфиденциальности в пределах манифеста \\
Сеть & не доверена; сообщения подписаны, повтор отсекается эпохой и номерами \\
\bottomrule
\end{tabular}
\end{center}

\section{Что гарантируется}

При соблюдении допущений уровня:

\begin{itemize}
\item ни одного эффекта на среде без совершённой и актуальной авторизации (I1);
\item ни одного продвижения основной ветки без финального одобрения и аттестаций, адресованных ровно к продвигаемому кандидату и текущему контракту (I3);
\item полномочия и условия приёмки задачи не превышают и не ослабляют выданного (I2);
\item оркестратор не выбирает ревьювера и не интегрирует в одиночку (I4, I5);
\item остановка доступна одной честной кор-ноде (I8);
\item каждое обращение к среде через протокол имеет проверяемое происхождение (I10).
\end{itemize}

\section{Чего не гарантируется}

\begin{warnbox}[Явные ограничения]
\begin{itemize}
\item Правильность суждения модели. Ревью независимо и подотчётно, но не доказано верным.
\item Предотвращение компрометации среды. Она обнаруживается по расхождению хешей и receipts, а при наличии верификатора по расхождению тестов, но не предотвращается.
\item Конфиденциальность кода от контрибьюторов и их провайдеров.
\item Устойчивость к побегу из санбокса через уязвимость ядра. Санбоксы не разделяют записываемых каталогов, а всё ценное вынесено за операции раннера, чтобы побег давал как можно меньше.
\item Единственность истории на L1 при злонамеренном секвенсере.
\item Полнота наблюдения за тем, что недоверенная нода делает вне протокола.
\item Живучесть при злонамеренном большинстве кор-нод.
\item Справедливость экономики.
\end{itemize}
\end{warnbox}

\section{Сценарий: злонамеренный оркестратор}

Оркестратор получил роль и раздаёт вредные задачи или занимается вандализмом. Четыре слоя:

\begin{enumerate}
\item У него нет рук. Среда исполняет только команды рабочих нод по назначениям с сертификатом кворума. Сам оркестратор действует лишь записями.
\item Каждая запись проходит валидаторов. Он не выбирает ревьювера, не интегрирует без чужого ревью и аттестации тестов с третьей ноды, не превышает лимиты, не выходит за capability, не трогает через воркеров защищённые пути.
\item Вредная по смыслу задача, прошедшая правила, исполняется и ревьюится посторонними нодами, а интеграцию подписывают другие кор-ноды, которые видят поток записей. Любая из них ставит систему на паузу одной подписью, после чего кворум сменяет оркестратора и исключает его из состава.
\item Ущерб обратим. Хеши в журнале и git позволяют откатить состояние к последнему вестированному, начисления вандала не вестятся, его аттестации отрицательны.
\end{enumerate}

Остаточный риск: правильно исполненная неправильная задача, интегрированная до того, как её заметили. Он ограничен областью capability, бюджетом, вестингом и остановкой, но не устранён.

\section{Сценарий: злонамеренный воркер}

Воркер получил задачу и пытается навредить. Слои: санбокс ограничивает эффекты областью задачи, без сети и без секретов; правило write scope отвергает ревизию, вышедшую за область; случайный ревьювер другого домена на другой модели видит диф; тесты на кандидате гоняет третья нода; вестинг задерживает начисление; аттестация репутации снижает вес ноды и увеличивает глубину проверки её следующих задач. Что воркер может сделать наверняка: сжечь свою ёмкость впустую и получить отрицательную аттестацию.

\section{Сценарий: скомпрометированная среда}

Среда может показать ревьюверу один диф, а интегрировать другой; вернуть поддельный результат проверок; выполнить другую команду; утечь код. Протокол это не предотвращает. Важно и то, чего он не обещает: receipt есть подписанное свидетельство самой среды, а не независимо установленный факт, и если среда скомпрометирована, её подпись не доказывает правдивость отчёта. Обнаружение поэтому условно: оно требует независимого наблюдения, верификатора на другой машине, который повторяет критические проверки, или собственных git-объектов у ревьювера, по которым можно сверить диф с тем, что записано. Журнал пинит хеши коммитов, одобрений и аттестаций, receipts несут дайджесты команд и состояний, и это делает расхождение выявимым тем, у кого есть с чем сравнить. Аппаратная аттестация исполнения могла бы перевести это из обнаружимого в предотвращаемое, но на macOS её нет, и первое развёртывание честно живёт с доверенной средой.

\section{Числа}

Пусть доля злонамеренных рабочих нод равна $p$, честный ревьювер обнаруживает злонамеренную ревизию с вероятностью $d$, ревьюверов $k$. Злонамеренная ревизия проходит ревью, если каждый ревьювер либо злонамерен, либо честен и пропустил:
\[
q = p + (1 - p)(1 - d), \qquad P = q^{k}.
\]

\begin{center}
\begin{tabular}{cc ccc}
\toprule
$p$ & $d$ & $k=1$ & $k=2$ & $k=3$ \\
\midrule
0.1 & 0.7 & 0.370 & 0.137 & 0.051 \\
0.1 & 0.9 & 0.190 & 0.036 & 0.007 \\
0.3 & 0.7 & 0.510 & 0.260 & 0.133 \\
0.3 & 0.9 & 0.370 & 0.137 & 0.051 \\
\bottomrule
\end{tabular}
\end{center}

Два вывода. Первый: вероятность обнаружения $d$ влияет сильнее, чем доля злонамеренных $p$, и $d$ есть свойство модели и протокола ревью, которое протокол не улучшает, а умножает: независимыми ревьюверами, тестами, верификатором, вестингом и репутацией. Второй: формула предполагает независимость ошибок, а три оси делают ошибки лишь менее коррелированными. Таблица есть иллюстрация, а не оценка безопасности; реальная величина измеряется на данных проекта.

\section{Обнаруживаемое и предотвращаемое}

\begin{center}
\small
\begin{tabular}{p{6.6cm} p{6.6cm}}
\toprule
Предотвращается & Обнаруживается, но не предотвращается \\
\midrule
эффекты без авторизации & компрометация среды \\
интеграция непроверенного состояния & семантически вредная задача внутри своей области \\
превышение полномочий & пропуск дефекта ревьювером \\
выбор ревьювера оркестратором & \\
повтор и откат записей & \\
блокировка остановки & \\
\bottomrule
\end{tabular}
\end{center}

Правый столбец есть честная цена того, что суждение вынесено за консенсус. Левый столбец есть то, чего не даёт ни одна из существующих систем параллельных агентов.
